% =========================================================
% Proof: Stern--Brocot Interval Refinement
% =========================================================

\newpage
\phantomsection
\label{prf:stern-brocot-refinement}

\begin{remark*}[Return]
\hyperref[thm:stern-brocot-refinement]{Return to Theorem~\ref*{thm:stern-brocot-refinement}.}
\end{remark*}

\begin{theorem*}[Stern--Brocot interval refinement]
Let $x > 0$.  Starting from $a/b < x < c/d$ with $b,d > 0$, repeatedly
insert the mediant $(a+c)/(b+d)$ and replace whichever endpoint keeps $x$
in the interval.  This generates a nested sequence of rational intervals
each containing $x$.
\end{theorem*}

\begin{proof}[Proof (Professional Standard)]
We show by induction that at every step $n$, the interval
$I_n = [a_n/b_n, c_n/d_n]$ satisfies $a_n/b_n < x < c_n/d_n$.

\textbf{Base case.} $a_0/b_0 = a/b < x < c/d = c_0/d_0$ by hypothesis.

\textbf{Inductive step.} Suppose $a_n/b_n < x < c_n/d_n$.  The mediant
$m_n = (a_n + c_n)/(b_n + d_n)$ satisfies $a_n/b_n < m_n < c_n/d_n$ by the
Mediant Inequality.  Compare $x$ to $m_n$:
\begin{itemize}
  \item If $x < m_n$: set $I_{n+1} = [a_n/b_n, m_n]$.  Then
        $a_n/b_n < x < m_n = c_{n+1}/d_{n+1}$.
  \item If $x > m_n$: set $I_{n+1} = [m_n, c_n/d_n]$.  Then
        $m_n = a_{n+1}/b_{n+1} < x < c_n/d_n$.
\end{itemize}
(Since $x$ is irrational in the case of an infinite path, $x = m_n$
cannot occur.)  In both cases $I_{n+1} \subsetneq I_n$ and $x \in I_{n+1}$.

By induction, $x \in I_n$ for all $n$ and $I_0 \supseteq I_1 \supseteq \cdots$.
\end{proof}

\begin{proof}[Proof (Detailed Learning)]
\textbf{Step 1: Induction statement.}
We claim: for all $n \in \mathbb{N}$, $a_n/b_n < x < c_n/d_n$.

\textbf{Step 2: Base case ($n = 0$).}
$a_0/b_0 = a/b < x < c/d = c_0/d_0$ by the hypothesis of the theorem.

\textbf{Step 3: Inductive step.}
Assume $a_n/b_n < x < c_n/d_n$.

Form the mediant $m_n = (a_n + c_n)/(b_n + d_n)$.  By the
\hyperref[thm:mediant-inequality]{Mediant Inequality}:
\[
\frac{a_n}{b_n} < m_n < \frac{c_n}{d_n}.
\]

Compare $x$ to $m_n$:
\begin{itemize}
  \item \textbf{Case $x < m_n$:} Define $a_{n+1}/b_{n+1} = a_n/b_n$ and
        $c_{n+1}/d_{n+1} = m_n$.  Then $a_n/b_n < x < m_n$, so
        $a_{n+1}/b_{n+1} < x < c_{n+1}/d_{n+1}$.  Also
        $I_{n+1} = [a_n/b_n, m_n] \subset [a_n/b_n, c_n/d_n] = I_n$.
  \item \textbf{Case $x > m_n$:} Define $a_{n+1}/b_{n+1} = m_n$ and
        $c_{n+1}/d_{n+1} = c_n/d_n$.  Then $m_n < x < c_n/d_n$, so
        $a_{n+1}/b_{n+1} < x < c_{n+1}/d_{n+1}$.  Also
        $I_{n+1} = [m_n, c_n/d_n] \subset [a_n/b_n, c_n/d_n] = I_n$.
\end{itemize}
In both cases the inductive claim holds at $n+1$ and $I_{n+1} \subsetneq I_n$.

\textbf{Step 4: Conclude.}
By induction, $x \in I_n$ for all $n$ and $(I_n)$ is a nested sequence of
rational intervals.
\end{proof}

\begin{remark*}[Proof structure]
An induction on steps, with the Mediant Inequality as the key tool at each
level.  The containment $x \in I_{n+1}$ follows by a two-case split on
whether $x < m_n$ or $x > m_n$; in both cases $x$ is retained in the
updated interval.
\end{remark*}

\begin{remark*}[Dependencies]
\hyperref[thm:mediant-inequality]{Mediant inequality},
\hyperref[def:mediant]{Mediant},
order properties of $\mathbb{R}$.
\end{remark*}
